% AA-22 | Algebra abstrakcyjna | Iteracja XI
\chapter[Twierdzenie Artina i odpowiedniość Galois]{Twierdzenie Artina i podstawowe twierdzenie teorii Galois}
\label{aa:ch:galois-correspondence}

\section{Od automorfizmów do ciał elementów stałych}
W Rozdziale~\ref{aa:ch:embeddings-normal-separable} zobaczyliśmy,
że skończone rozszerzenie Galois $L/K$ ma dokładnie
$[L:K]$ automorfizmów pozostawiających elementy $K$ bez zmian.
Teraz odwracamy pytanie: dla zbioru automorfizmów chcemy
odtworzyć podciało, które jest przez nie ustalone, oraz
ustalić, czy \emph{każde} ciało pośrednie można tak opisać.

\begin{definicja}[Ciało elementów stałych]
\label{aa:def:fixed-field}
Niech $L$ będzie ciałem, a $H\le\operatorname{Aut}(L)$
podgrupą grupy automorfizmów ciała. Definiujemy
\emph{ciało elementów stałych} podgrupy $H$ przez
\[
L^H=\{x\in L:\sigma(x)=x\text{ dla każdego }\sigma\in H\}.
\]
Jeśli $L/K$ jest rozszerzeniem ciał i
$H\le\operatorname{Gal}(L/K)$, to $K\subseteq L^H\subseteq L$.
\end{definicja}

\begin{lemat}[Podstawowe własności ciała elementów stałych]
\label{aa:lem:fixed-field}
Niech $L$ będzie ciałem i niech $H\le\operatorname{Aut}(L)$. Wtedy $L^H$ jest
podciałem $L$. Jeśli $H_1,H_2\le\operatorname{Aut}(L)$ i
$H_1\subseteq H_2$, to $L^{H_2}\subseteq L^{H_1}$.
Ponadto dla
$\tau\in\operatorname{Aut}(L)$ zachodzi
$\tau(L^H)=L^{\tau H\tau^{-1}}$.
\end{lemat}
\begin{proof}
Każdy automorfizm zachowuje $0$ i $1$, więc $0,1\in L^H$.
Jeśli $x,y\in L^H$, to dla każdego $\sigma\in H$ mamy
$\sigma(x-y)=x-y$, $\sigma(xy)=xy$, a dla $x\ne0$
także $\sigma(x^{-1})=\sigma(x)^{-1}=x^{-1}$.
Dowodzi to, że $L^H$ jest podciałem. Włączenie dla
$H_1\subseteq H_2$ jest bezpośrednie: warunek
niezmienniczości względem większej grupy jest silniejszy.
Wreszcie $x\in L^H$ jest równoważne równości
$(\tau\sigma\tau^{-1})(\tau(x))=\tau(x)$ dla
wszystkich $\sigma\in H$. Ponieważ $\tau$ jest bijekcją,
daje to ostatnią tożsamość.
\end{proof}

\begin{przyklad}[Pierwsza para: ciało i podgrupa]
\label{aa:prz:fixed-quadratic}
W $L=\mathbb Q(\sqrt2)$ niech $H$ będzie grupą
złożoną z identyczności i automorfizmu
$\tau(a+b\sqrt2)=a-b\sqrt2$. Wtedy
$\tau(a+b\sqrt2)=a+b\sqrt2$ wtedy i tylko wtedy,
gdy $2b\sqrt2=0$, czyli $b=0$.
Dlatego $L^H=\mathbb Q$. Natomiast dla podgrupy
trywialnej $\{\operatorname{id}_L\}$ mamy
$L^{\{\operatorname{id}_L\}}=L$.
Widać zatem, że \emph{większa podgrupa odpowiada
mniejszemu ciału elementów stałych}.
\end{przyklad}

\section{Liniowa niezależność różnych zanurzeń}
Poniższy lemat jest zasadniczym narzędziem algebraicznym.
Nie korzysta ani z twierdzenia o odpowiedniości Galois,
ani z wcześniej znanego wzoru na stopień ciała elementów stałych.

\begin{lemat}[Niezależność homomorfizmów ciał]
\label{aa:lem:artin-independence}
Niech $L,M$ będą ciałami, niech $r\ge1$ będzie liczbą
całkowitą i niech
$\sigma_1,\ldots,\sigma_r:L\to M$ będą parami różnymi
homomorfizmami ciał. Jeśli $c_1,\ldots,c_r\in M$ spełniają
\[
\sum_{i=1}^r c_i\sigma_i(x)=0\qquad\text{dla każdego }x\in L,
\]
to wszystkie $c_i$ są zerowe. Innymi słowy, homomorfizmy
te są liniowo niezależne nad $M$ jako funkcje $L\to M$.
\end{lemat}
\begin{proof}
Przypuśćmy, że istnieje nietrywialna relacja, i wybierzmy
taką relację z najmniejszą możliwą liczbą niezerowych
współczynników. Po pominięciu wyrazów zerowych mamy
$\sum_{i=1}^s c_i\sigma_i(x)=0$ dla każdego $x\in L$,
gdzie wszystkie $c_i\ne0$. Nie może być $s=1$, bo
podstawienie $x=1$ dałoby $c_1=0$.
Ponieważ $\sigma_1\ne\sigma_s$, istnieje $a\in L$ takie,
że $\sigma_1(a)\ne\sigma_s(a)$.
W relacji podstawiamy $ax$ w miejsce $x$, a następnie
odejmujemy od niej pierwotną relację pomnożoną przez
$\sigma_s(a)$. Dla każdego $x\in L$ dostajemy
\[
\sum_{i=1}^{s-1}c_i(\sigma_i(a)-\sigma_s(a))\sigma_i(x)=0.
\]
Jest to relacja o mniejszej liczbie wyrazów, lecz
współczynnik przy $\sigma_1$ jest niezerowy.
Sprzeczność z minimalnością kończy dowód.
\end{proof}

\section{Twierdzenie Artina o ciele elementów stałych}
\begin{twierdzenie}[Twierdzenie Artina]
\label{aa:tw:artin-fixed-field}
Niech $L$ będzie ciałem, a $G$ skończoną podgrupą
grupy $\operatorname{Aut}(L)$. Niech $F=L^G$.
Wtedy $L/F$ jest skończonym rozszerzeniem Galois,
\[
[L:F]=|G|,\qquad \operatorname{Gal}(L/F)=G.
\]
Nie zakładamy z góry, że $L$ ma skończony stopień
nad $F$.
\end{twierdzenie}
\begin{proof}
Oznaczmy $m=|G|$. Najpierw pokażemy $[L:F]\le m$.
Wybierzmy dowolne $m+1$ elementów $x_1,\ldots,x_{m+1}\in L$.
Rozważmy układ $m$ jednorodnych równań o niewiadomych
$a_1,\ldots,a_{m+1}\in L$:
\[
\sum_{j=1}^{m+1}a_j\sigma(x_j)=0
\qquad(\sigma\in G).
\]
Ponieważ liczba niewiadomych jest większa od liczby
równań, układ ma niezerowe rozwiązanie.
Wybierzmy spośród takich rozwiązań jedno z najmniejszą
możliwą liczbą niezerowych współczynników.
Dzieląc wszystkie współczynniki przez dowolny wybrany
niezerowy współczynnik, możemy przyjąć $a_{j_0}=1$
dla pewnego $j_0$.

Dla $\tau\in G$ zastosujmy $\tau$ do równania
odpowiadającego $\tau^{-1}\sigma$.
Ponieważ $\tau((\tau^{-1}\sigma)(x_j))=\sigma(x_j)$,
otrzymujemy
$\sum_j\tau(a_j)\sigma(x_j)=0$ dla każdego $\sigma\in G$.
Tak więc rodzina $(\tau(a_j))_j$ także jest rozwiązaniem.
Różnica $(\tau(a_j)-a_j)_j$ jest rozwiązaniem,
a jej współczynnik przy $j_0$ jest równy
$\tau(1)-1=0$. Jeżeli różnica byłaby niezerowa,
miałaby mniej niezerowych współczynników niż wybrane
minimalne rozwiązanie, co jest niemożliwe.
Stąd $\tau(a_j)=a_j$ dla każdego $j$ i $\tau\in G$,
a więc wszystkie $a_j$ należą do $F$.
Dla $\sigma=\operatorname{id}_L$ dostajemy
$\sum_j a_jx_j=0$, nietrywialną zależność liniową
nad $F$: co najmniej jeden współczynnik $a_j$ jest
niezerowy, a wszystkie należą do $F$.
Ponieważ dowolne $m+1$ elementów jest
zależne nad $F$, mamy $[L:F]\le m$.

Oznaczmy teraz $n=[L:F]$, który jest skończony,
i wybierzmy bazę $b_1,\ldots,b_n$ przestrzeni $L$ nad $F$.
Każdy $\sigma\in G$ jest $F$-liniowy, a więc wyznaczony
przez wartości $\sigma(b_1),\ldots,\sigma(b_n)$.
Gdyby $m>n$, to $m$ wektorów
$(\sigma(b_1),\ldots,\sigma(b_n))\in L^n$ byłoby
liniowo zależnych nad $L$. Istniałyby zatem
współczynniki $c_\sigma\in L$, nie wszystkie zerowe,
takie, że $\sum_{\sigma\in G}c_\sigma\sigma(b_j)=0$
dla każdego $j$. Dzięki $F$-liniowości tych automorfizmów
równość przeniosłaby się na wszystkie $x\in L$, wbrew
Lematowi~\ref{aa:lem:artin-independence}.
Zatem $m\le n$, co wraz z pierwszą nierównością daje
$[L:F]=m$.

Dla $x\in L$ rozważmy jego orbitę $Gx=\{\sigma(x):\sigma\in G\}$.
Wielomian
\[
p_x(X)=\prod_{y\in Gx}(X-y)\in L[X]
\]
ma parami różne pierwiastki. Grupa $G$ permutuje orbitę,
więc każdy współczynnik $p_x$ jest ustalany przez
każdy element $G$; dlatego $p_x\in F[X]$.
Wielomian minimalny $m_{x,F}$ dzieli $p_x$ w $F[X]$,
ponieważ $p_x(x)=0$. Wielomian $p_x$ ma w $L$ same
pierwiastki proste, więc również $m_{x,F}$ ma tam same
pierwiastki proste i rozkłada się na czynniki liniowe.
Ponieważ dotyczy to każdego $x\in L$, rozszerzenie
$L/F$ jest rozdzielcze. Jest też normalne na mocy
kryterium rozkładu wszystkich wielomianów minimalnych
z Twierdzenia~\ref{aa:tw:normal-embeddings}.
Jest skończone, zatem Galois.
Wreszcie $G\subseteq\operatorname{Gal}(L/F)$, ponieważ
każdy automorfizm z $G$ ustala wszystkie elementy $F=L^G$.
Z Twierdzenia~\ref{aa:tw:galois-group-degree} wynika,
że grupa po prawej ma $[L:F]=m=|G|$ elementów.
Obie grupy są więc równe.
\end{proof}

\begin{wniosek}[Odzyskiwanie grupy z jej ciała elementów stałych]
\label{aa:wn:artin-recover-group}
Jeśli $L/K$ jest skończonym rozszerzeniem Galois,
a $H\le\operatorname{Gal}(L/K)$, to
$\operatorname{Gal}(L/L^H)=H$ oraz $[L:L^H]=|H|$.
\end{wniosek}
\begin{proof}
Zastosujmy Twierdzenie~\ref{aa:tw:artin-fixed-field}
do skończonej grupy $H$ automorfizmów ciała $L$.
\end{proof}

\section{Podstawowe twierdzenie teorii Galois}
\begin{twierdzenie}[Podstawowe twierdzenie teorii Galois]
\label{aa:tw:fundamental-galois}
Niech $L/K$ będzie \emph{skończonym rozszerzeniem Galois}
i niech $G=\operatorname{Gal}(L/K)$. Między zbiorem
wszystkich ciał pośrednich $E$, $K\subseteq E\subseteq L$,
a zbiorem wszystkich podgrup $H\le G$ istnieją
wzajemnie odwrotne odpowiedniości
\[
E\longmapsto\operatorname{Gal}(L/E),\qquad
H\longmapsto L^H.
\]
Odwracają one inkluzje. Jeśli $E=L^H$, to
\[
[L:E]=|H|,\qquad [E:K]=[G:H].
\]
Ponadto $E/K$ jest normalne (a zatem jest rozszerzeniem
Galois) wtedy i tylko wtedy, gdy $H$ jest podgrupą normalną
w $G$. W tym przypadku ograniczenie automorfizmów daje
kanoniczny izomorfizm
$\operatorname{Gal}(E/K)\cong G/H$.
\end{twierdzenie}
\begin{proof}
\emph{Krok 1: odpowiedniości są odwrotne.}
Ustalmy ciało algebraicznie domknięte $\Omega$ zawierające $L$,
którego istnienie gwarantuje Twierdzenie~\ref{aa:tw:algebraic-closure-exists}.
Niech $E$ będzie ciałem pośrednim. Rozszerzenie $L/E$
jest skończone, ponieważ na mocy twierdzenia o wieży
(Twierdzenie~\ref{aa:tw:field-tower}) zachodzi
$[L:K]=[L:E][E:K]<\infty$.
Ponadto jest rozdzielcze: dla każdego $x\in L$
wielomian minimalny $x$ nad $E$ dzieli wielomian
minimalny nad $K$ w $E[X]$; ten drugi nie ma
pierwiastków wielokrotnych, więc pierwszy także nie.
Rozszerzenie $L/E$ jest również normalne. Z
Twierdzenia~\ref{aa:tw:normal-embeddings} wynika,
że istnieją wielomiany $f_1,\ldots,f_r\in K[X]$,
których ciałem rozkładu nad $K$ jest $L$.
Oznaczmy przez $S$ zbiór wszystkich ich pierwiastków
w $L$. Wtedy $L=K(S)$; ponieważ
$K\subseteq E\subseteq L$, mamy również $L=E(S)$.
Zatem te same wielomiany należą do $E[X]$, rozkładają
się w $L$ i generują $L$ nad $E$.
Stąd $L/E$ jest ciałem rozkładu tej rodziny nad $E$,
a zatem rozszerzeniem normalnym i rozdzielczym,
czyli Galois.

Niech $H=\operatorname{Gal}(L/E)$. Wszystkie elementy
$H$ ustalają $E$, zatem $E\subseteq L^H$.
Z Twierdzenia~\ref{aa:tw:galois-group-degree} i Wniosku
\ref{aa:wn:artin-recover-group} dostajemy
\[
[L:E]=|H|=[L:L^H].
\]
Ponieważ $E\subseteq L^H$, twierdzenie o wieży
wymusza $[L^H:E]=1$, a więc $E=L^H$.
Odwrotnie, dla każdej podgrupy $H\le G$ z Wniosku
\ref{aa:wn:artin-recover-group} wynika bezpośrednio
$\operatorname{Gal}(L/L^H)=H$.

\emph{Krok 2: inkluzje i stopnie.} Jeśli
$E_1\subseteq E_2$, to każdy automorfizm ustalający
$E_2$ ustala też $E_1$, czyli
$\operatorname{Gal}(L/E_2)\subseteq\operatorname{Gal}(L/E_1)$.
Dla podgrup $H_1\subseteq H_2$
Lemat~\ref{aa:lem:fixed-field} daje
$L^{H_2}\subseteq L^{H_1}$, więc także druga
odpowiedniość odwraca inkluzje. Dla $E=L^H$
Twierdzenie Artina daje $[L:E]=|H|$. Z twierdzenia
o wieży i równości $[L:K]=|G|$ wynika
$[E:K]=|G|/|H|=[G:H]$.

\emph{Krok 3: normalność.} Dla $\sigma\in G$ mamy
$\sigma(E)=L^{\sigma H\sigma^{-1}}$ na mocy
Lematu~\ref{aa:lem:fixed-field}.
Dzięki bijektywności odpowiedniości
$\sigma(E)=E$ wtedy i tylko wtedy, gdy
$\sigma H\sigma^{-1}=H$.
Zatem $H\triangleleft G$ dokładnie wtedy, gdy
$E$ jest zachowywane przez wszystkie automorfizmy
$\sigma\in G$; w obu kierunkach korzystamy z
jednoznaczności podgrupy odpowiadającej $E$.

Wykażmy, że to ostatnie jest równoważne normalności
$E/K$. Jeśli $E/K$ jest normalne, to dowolne
$K$-zanurzenie $E\hookrightarrow L$, w szczególności
ograniczenie $\sigma|_E$, ma obraz równy $E$:
wystarczy zastosować argument z dowodu
Twierdzenia~\ref{aa:tw:normal-embeddings} do ciała $E$.
Odwrotnie, załóżmy, że $\sigma(E)=E$ dla wszystkich
$\sigma\in G$. Dowolne $K$-zanurzenie
$\iota:E\hookrightarrow\Omega$ przedłuża się na
$K$-zanurzenie $\widetilde\iota:L\hookrightarrow\Omega$
na mocy Wniosku~\ref{aa:wn:extend-embedding-finite}.
Normalność $L/K$ gwarantuje, że
$\widetilde\iota\in G$. Zatem
$\iota(E)=\widetilde\iota(E)=E$.
Kryterium normalności przez zanurzenia daje
normalność $E/K$.

\emph{Krok 4: grupa ilorazowa.} Jeśli $H\triangleleft G$,
to $E/K$ jest normalne. Rozdzielczość przechodzi
na podrozszerzenia, więc $E/K$ jest Galois.
Ograniczenie
$\rho:G\to\operatorname{Gal}(E/K)$,
$\rho(\sigma)=\sigma|_E$, jest dobrze określonym
homomorfizmem, bo $\sigma(E)=E$. Jego jądrem są
dokładnie automorfizmy ustalające $E$, czyli $H$.
Jest on surjektywny: każdy $K$-automorfizm $\tau:E\to E$
przedłuża się na $K$-zanurzenie $L\to\Omega$,
które dzięki normalności $L/K$ jest elementem $G$.
Twierdzenie o izomorfizmie grup daje
$G/H\cong\operatorname{Gal}(E/K)$.
\end{proof}

\section{Przykłady odpowiedniości Galois}
\begin{przyklad}[Rozszerzenie dwukwadratowe]
\label{aa:prz:biquadratic-galois-lattice}
Niech $L=\mathbb Q(\sqrt2,\sqrt3)$. W
Przykładzie~\ref{aa:prz:biquadratic-extension} wykazaliśmy
$[L:\mathbb Q]=4$. Ciało $L$ jest ciałem rozkładu
$(X^2-2)(X^2-3)$, a rozszerzenie jest rozdzielcze,
ponieważ $\operatorname{char}\mathbb Q=0$.
Zatem $L/\mathbb Q$ jest rozszerzeniem Galois.

Niech $s$ zmienia znak $\sqrt2$ i zachowuje $\sqrt3$,
a $t$ zachowuje $\sqrt2$ i zmienia znak $\sqrt3$.
Istnienie tych automorfizmów wymaga uzasadnienia, a nie
samego wskazania możliwych obrazów generatorów.
Najpierw $\mathbb Q(\sqrt2)$ ma dwa zanurzenia nad
$\mathbb Q$, wysyłające $\sqrt2$ odpowiednio na
$\sqrt2$ i $-\sqrt2$ (Przykład~\ref{aa:prz:quadratic-conjugation}).
Wielomian minimalny $\sqrt3$ nad $\mathbb Q(\sqrt2)$
ma stopień dwa: gdyby $\sqrt3$ należał do
$\mathbb Q(\sqrt2)$, stopień rozszerzenia
$L/\mathbb Q$ byłby równy dwa, a wynosi cztery.
Jest to zatem $X^2-3$. Po przeniesieniu współczynników
przez każde z dwóch zanurzeń pozostaje on tym samym
wielomianem i ma dwa pierwiastki $\sqrt3,-\sqrt3$
w $L$. Lemat~\ref{aa:lem:extend-embedding-simple}
przedłuża każde zanurzenie na dwa sposoby. Powstałe
cztery zanurzenia $L\to L$ są automorfizmami:
każde jest $\mathbb Q$-liniowe i różnowartościowe,
a $L$ ma skończony wymiar nad $\mathbb Q$, więc
jest również surjektywne. W szczególności istnieją $s$ i $t$
o podanych wartościach na generatorach.
Mamy $s^2=t^2=1$, $st=ts$, a cztery automorfizmy to
$1,s,t,st$. Stąd
$G\cong C_2\times C_2$. Podgrupy rzędu dwa to
$\langle s\rangle$, $\langle t\rangle$ i
$\langle st\rangle$; ich ciała elementów stałych
mają stopień dwa nad $\mathbb Q$ i wynoszą odpowiednio
\[
L^{\langle s\rangle}=\mathbb Q(\sqrt3),\quad
L^{\langle t\rangle}=\mathbb Q(\sqrt2),\quad
L^{\langle st\rangle}=\mathbb Q(\sqrt6).
\]
Ostatnia równość wynika z tego, że $st$ zmienia
oba znaki, a więc zachowuje $\sqrt6$, oraz ze
zgodności stopni obu ciał nad $\mathbb Q$.
Wszystkie podgrupy $G$ są normalne, dlatego każde
ciało pośrednie jest normalne nad $\mathbb Q$.
\end{przyklad}

\begin{przyklad}[Grupa symetryczna i nienormalne ciało pośrednie]
\label{aa:prz:cubic-galois-s3}
Niech $\alpha=\sqrt[3]2\in\mathbb R$,
$\omega\in\mathbb C$ będzie pierwiastkiem trzeciego
stopnia z jedynki różnym od $1$, a
$L=\mathbb Q(\alpha,\omega)$. Z
Przykładu~\ref{aa:prz:cubic-splitting} wiemy,
że $L$ jest ciałem rozkładu $X^3-2$ i
$[L:\mathbb Q]=6$. Wielomian jest rozdzielczy
w charakterystyce zero, zatem grupa $G=\operatorname{Gal}(L/\mathbb Q)$
ma sześć elementów.

Definiujemy dwa automorfizmy na generatorach:
\[
\sigma(\alpha)=\omega\alpha,\quad\sigma(\omega)=\omega,
\qquad
\tau(\alpha)=\alpha,\quad\tau(\omega)=\omega^2.
\]
Aby uzasadnić istnienie $\sigma$, zauważmy, że
$[\mathbb Q(\omega):\mathbb Q]=2$, a
$[L:\mathbb Q]=6$, zatem $[L:\mathbb Q(\omega)]=3$.
Wielomian $X^3-2$ jest więc wielomianem minimalnym
$\alpha$ nad $\mathbb Q(\omega)$ i ma w $L$ również
pierwiastek $\omega\alpha$.
Lemat~\ref{aa:lem:extend-embedding-simple} daje zatem
$\mathbb Q(\omega)$-zanurzenie $L\to L$ o wskazanych
wartościach. Jego obraz zawiera $\omega$ oraz
$\omega\alpha$, więc także
$\alpha=(\omega\alpha)/\omega$, a zatem całe $L$.
Jest to automorfizm $\sigma$. Sprzężenie zespolone
zachowuje $L$, ponieważ $\alpha\in\mathbb R$ oraz
$\overline{\omega}=\omega^2\in L$. Jego ograniczenie
do $L$ jest drugim $\mathbb Q$-automorfizmem $\tau$. Mamy
$\sigma^3=\tau^2=1$ oraz
$\tau\sigma\tau^{-1}=\sigma^{-1}$.
Elementy $1,\sigma,\sigma^2,\tau,\sigma\tau,\sigma^2\tau$
są parami różne: trzy pierwsze ustalają $\omega$,
a trzy ostatnie wysyłają go na $\omega^2$;
w każdej trójce obrazy $\alpha$ są różne.
Ponieważ $|G|=6$, elementy te tworzą całe $G$.
Działanie $G$ na trzech pierwiastkach
$\alpha,\omega\alpha,\omega^2\alpha$ daje homomorfizm
do $S_3$. Jest on różnowartościowy, bo automorfizm
ustalający wszystkie trzy pierwiastki ustala zarówno
$\alpha$, jak i $\omega=(\omega\alpha)/\alpha$,
a więc całe $L$. Obie grupy mają rząd sześć,
zatem $G\cong S_3$.

Podgrupa $\langle\sigma\rangle$ ma rząd trzy,
więc jej ciało elementów stałych ma stopień dwa
nad $\mathbb Q$. Ponieważ $\sigma(\omega)=\omega$,
mamy $L^{\langle\sigma\rangle}=\mathbb Q(\omega)$.
Podgrupa $\langle\tau\rangle$ ma rząd dwa,
a jej ciało elementów stałych to $\mathbb Q(\alpha)$:
$\tau$ zachowuje $\alpha$, a oba ciała mają stopień
trzy nad $\mathbb Q$. Podgrupa $\langle\tau\rangle$
nie jest normalna w $S_3$, więc
$\mathbb Q(\alpha)/\mathbb Q$ nie jest normalne.
To pokazuje na konkretnym przykładzie, dlaczego
w podstawowym twierdzeniu teorii Galois nie można
żądać normalności każdego ciała pośredniego.
\end{przyklad}

\begin{przyklad}[Podgrupy i ciała pośrednie ciała skończonego]
\label{aa:prz:finite-fields-galois-lattice}
Niech $p$ będzie liczbą pierwszą,
$L=\mathbb F_{p^{12}}$ i $K=\mathbb F_p$
(przy czym $K$ utożsamiamy z pierwszym podciałem $L$).
Grupa Galois rozszerzenia $L/K$ jest cykliczna rzędu $12$, generowana
przez $\operatorname{Fr}_p:x\mapsto x^p$.
Dla każdego dzielnika $d\mid12$ istnieje dokładnie
jedna podgrupa indeksu $d$, generowana przez
$\operatorname{Fr}_p^{d}$. Jej ciało elementów stałych
ma stopień $d$ nad $K$ i dlatego $p^d$ elementów.
Zatem dla każdego $d\in\{1,2,3,4,6,12\}$
odpowiedniość Galois daje dokładnie jedno takie
ciało pośrednie. Z twierdzenia o jednoznaczności ciał
skończonych wynika, że jest ono izomorficzne z
$\mathbb F_{p^d}$.
Wszystkie te rozszerzenia są normalne, ponieważ
każda podgrupa grupy cyklicznej jest normalna.
\end{przyklad}

\begin{uwaga}[Granice obecnego etapu]
\label{aa:uw:galois-future}
Podstawowe twierdzenie teorii Galois klasyfikuje ciała
pośrednie skończonych rozszerzeń Galois poprzez podgrupy.
Nie rozstrzyga jeszcze, kiedy dane równanie wielomianowe
jest rozwiązywalne przez pierwiastniki. W następnej części
trzeba zestawić tę odpowiedniość z szeregiem podnormalnym
i rozwiązalnością grup z Rozdziału~\ref{aa:ch:szeregi}.
Osobnego opracowania wymagają również norma, ślad,
rozszerzenia cykliczne i odpowiednie zastosowania teorii Galois.
\end{uwaga}

\begin{uwaga}[Dlaczego odpowiedniość wymaga rozszerzenia Galois]
Niech $L=\mathbb Q(\sqrt[3]{2})$ i $K=\mathbb Q$.
Istotnie, każdy $K$-automorfizm $L$ musi wysłać
$\sqrt[3]{2}$ na pierwiastek jego wielomianu minimalnego
$X^3-2$. Ponieważ $L\subseteq\mathbb R$, jedynym takim
pierwiastkiem w $L$ jest $\sqrt[3]{2}$, więc automorfizm
jest identycznością. Grupa $\operatorname{Aut}_K(L)$
jest zatem trywialna, a jej ciało elementów stałych
to całe $L$, nie zaś $K$.
Nie wolno zatem przenosić równości
$L^{\operatorname{Gal}(L/K)}=K$ na dowolne rozszerzenia
skończone bez założenia normalności.
\end{uwaga}
